% Nouveau EGA V (cote 25): archive cote 25, site page 137 (PDF page 138).
% English translation of the checked French annotated master.
% AI-translated by GPT-6 Astra (OpenAI, ChatGPT Pro mode).
\scanpage{25/137}
\mspage{\hand{40}}

IV 21 Invertible sheaves and relative divisors on projective
and multiprojective bundles; linear systems of divisors.

1. Invertible sheaves on a projective bundle; application to
morphisms from one projective bundle to another.

2. Representability of $\underline{\mathrm{Div}}^{L}_{X/S}$; \var{\inserted[typed]{relative}}{relative}{insertion}{ega25_25_137-v0002} divisors on projective and
multiprojective bundles.

3. Linear systems of divisors$_{\mathrm{x}}$ and morphisms into projective
bundles.

4. Linear systems of divisors$_{\mathrm{x}}$ and invertible Modules.

The results developed in this section and the following ones
are already partly global in nature and give, in particular, further
results on projective schemes, using \var{\xout{in an essential way}}{}{typewriter strike-out}{ega25_25_137-v0003} the
global constructions of Chap II, and even here certain results of
Chap III, in addition to the purely local results of this Chapter IV.

One of the aims of this section is to develop the language of
``linear systems of divisors'', linked on the one hand to the classification of
morphisms into a projective bundle, and on the other to that of invertible Modules
on a given relative prescheme. Let us observe here that the
truly natural ``parameter schemes'' for linear systems of
divisors are the \uline{Brauer-Séveri} schemes, which generalize projective
bundles and can be defined, for example, as bundles that become
isomorphic to a projective bundle after a surjective étale extension of
the base. Since their study uses descent theory, developed
in Chap V, and since moreover their classification is equivalent to that of
\var{\xout{\unsure{principal homogeneous bundles}}}{}{typewriter strike-out}{ega25_25_137-v0005} torsors with group the projective group, we
refer to Chap VI$_{\mathrm{x}}$ of our learned work for the study of these schemes and their links with the notion of a
linear system of divisors.

\mspage{\hand{137}}
\tnote{25-137-1}{The canceled fragments lesxf and fibr have no separate English lexical content and are recorded as omitted French-only cancellations. The source spellings systêmes and desur, and the split inver/sibles, have not been imitated with English spelling errors. The lowered x characters following divisors and VI remain exactly as in the French master.}
